\section{模型评价与推广}

\subsection{模型的优点}
\textbf{（1）预测模型由数据规律驱动，可解释性强。} 本文并未直接套用通用时间序列方法，
而是先对全年数据作统计分析，发现负载的星期型双状态结构
（周日至周四约 124.0\,MWh，周五、周六约 78.3\,MWh），
据此选择“上周同星期实际值”作为预测基础，避免持续性预测在状态切换日失效；
光伏出力具有强季节性与日间持续性（自相关系数 0.964），支持使用近期历史数据。
模型变量与参数均有明确的实际含义。

\textbf{（2）关键参数具有理论依据而非经验试凑。} 5 倍电价惩罚使购电决策具有报童模型结构，
临界分位数 0.8 要求计划覆盖净负荷约 80\% 分位点，据此设置的安全余量在全年扫描中形成平坦平台
（$\gamma$ 在 $[1.03,1.07]$ 内总费用变化不足 1\%），分段验证结论一致，参数稳健、无过拟合。

\textbf{（3）模型分层递进，四个问题共用同一优化内核。} 从问题一的日循环线性规划，
到问题二“预测—余量—日前计划—日内纠偏—跨日滚动”，
再到问题三“预报标定—残差修正—双向调整—末端备用”与问题四的价格项替换，
模型结构始终一致、各层功能清晰，便于横向比较不同策略的经济性。

\textbf{（4）求解可靠、结果可复现且具有经济含义。} 日前计划与调整计划均为线性规划，可求得全局最优；
日内纠偏规则透明、物理可行。全链路数值校验在 334 天滚动仿真中零违规，
Python 与 MATLAB 双平台求解结果逐位一致；单次全年仿真在数分钟内完成，使参数大规模扫描寻优成为可能。
从结果看，储能设备年价值 439.2 万元；问题三把 5 倍紧急购电转化为 0.5 倍与 1.5 倍调整结算后，
全年费用较问题二下降 5.9\%，紧急购电量由 14.51 万\,kWh 降至 5.59 万\,kWh，与费率结构改善的机理吻合。

\subsection{模型的不足}
\textbf{（1）预测仍为点预测加比例余量。} 未建立负载与光伏的联合概率分布或场景集，
余量按全天统一设置、未按时段细化，存在固有代价（问题二全年弃电约 326 万\,kWh，
其中已付费却被迫丢弃的电量下限约 81 万\,kWh）。
采用分布鲁棒优化\cite{bertsimas2013}或分时段余量有望进一步压缩该损失。

\textbf{（2）日内调度采用贪心规则。} 储能纠偏按固定优先级执行，未在日内逐时段重新优化，
与模型预测控制（MPC）\cite{mayne2000}相比仍留有优化空间（估计 1\%$\sim$2\% 量级）。

\textbf{（3）电价不确定性与参数寻优方式尚不完善。} 问题四假设计划制定时当日电价已知，
电价未知时应引入价格场景集或不确定集，将模型扩展为两阶段随机规划；
余量系数、残差权重与备用水平均以全年数据寻优，虽已通过平台分析与分段验证排除过拟合，
严格意义上仍应采用滚动交叉验证。

\textbf{（4）目标函数未考虑储能循环损耗。} 模型以购电费用最小为唯一目标，
未计入充放电次数带来的寿命损耗；引入单位循环成本后储能周转量将相应下降。

\subsection{模型的推广}
本文的核心思想——\textbf{在不对称惩罚下用保守余量制定日前计划、用日内滚动调整校正预测偏差}——
不依赖具体场景，可推广到其他“日前承诺、实时考核”型决策问题。

\textbf{（1）同类电力系统与设备。} 风电场与光伏电站的日前出力申报、独立储能电站的峰谷套利、
园区综合能源系统调度、需求响应中可中断负荷的报价等，
均具有“计划量按合同结算、偏差按惩罚电价结算”的结构，可直接沿用本文两阶段框架，
只需替换考核系数（5 倍、1.5 倍、0.5 倍）与设备的容量—功率—效率参数。

\textbf{（2）数据与机制变化下的调整。} 预报发布更频繁时可加密调整时刻；
电价未知时可将固定电价替换为价格场景集，扩展为两阶段随机规划或鲁棒优化；
充放电效率随荷电状态变化时可将线性储能模型替换为分段线性模型。

\textbf{（3）方法层面的改进方向。} 预测环节可用 LSTM、Transformer 等机器学习模型或
概率预测方法\cite{hong2016}替代 MOS 线性回归；日内调度可由贪心规则升级为模型预测控制；
还可引入储能寿命模型，把循环损耗纳入目标函数。
